\documentclass[10pt,a4paper]{amsart}
\usepackage[margin=19mm]{geometry}
\usepackage{amsmath,amssymb,mathtools}
\usepackage[scr=rsfs,cal=euler]{mathalfa}
\usepackage{xfrac}
\usepackage[unicode,hidelinks,bookmarks=false]{hyperref}
\newcommand{\ie}{i.e.\ }
\newcommand{\cf}{cf.\ }
\newcommand{\resp}{respectively\ }
\newcommand{\Subsection}{\subsection}
\providecommand{\nobreakdash}{\nobreak\hbox{-}}
\setlength{\emergencystretch}{2em}
\multlinegap0pt
\usepackage{bm}
\usepackage{bbm}
\usepackage{dsfont}
\usepackage{xspace}
\usepackage{cancel}
\usepackage{mathtools}
\usepackage{amsthm}
\makeatletter
\@ifundefined{lemma}{\newtheorem{lemma}{Lemma}}{}
\@ifundefined{theorem}{\newtheorem{theorem}{Theorem}}{}
\makeatother
\DeclareMathOperator{\codim}{codim}
\title[Dichotomies for triangular systems on Hilbert spaces]{Dichotomies for triangular systems on Hilbert spaces}
\author{Davor Dragicevic, Kenneth J. Palmer, Boris Petkovic}
\date{Source: arXiv:2508.04320v1 (CC BY 4.0)}
\begin{document}
\maketitle

\noindent\emph{Source and licence.} Dichotomies for triangular systems on Hilbert spaces. Version arXiv:2508.04320v1 (CC BY 4.0), distributed under the Creative Commons Attribution 4.0 International licence, as recorded in the arXiv licence metadata for this exact version and shown on the abstract page https://arxiv.org/abs/2508.04320v1. Full rights notice: licenses/arxiv-2508.04320-CC-BY-4.0.txt.\par\smallskip

\noindent\textit{Editorial scope.} Counted unit: the Theorem whose source label is \texttt{515thm} in the article (source0.tex lines 1132--1135, source label \texttt{515thm}), delivered together with the article's own complete original proof of it (source0.tex lines 1137--1214, one \texttt{proof} environment of 78 lines). No mathematics is written, repaired or re-derived: statement and proof are the article's own text line for line. Required context retained uncounted: LDE (lines 269--274); ldedd (lines 276--281); iid (lines 293--295); Si (lines 357--359); thm7 (lines 703--711); ddd (lines 741--743). Every definition, display and label of those lines is retained unchanged. Omitted from this package and not redistributed: 1--268, 275--275, 282--292, 296--356, 360--702, 712--740, 744--1131, 1136--1136, 1215--1483. Nothing inside a retained excerpt is omitted or altered: every retained body is delivered exactly as the article's lines are written, so no omission receipt of any kind accompanies this package. Citations: the retained excerpts carry no citation command at all, so no bibliography entry of the article is transcribed and no reference is redistributed. Reference adaptation: none was needed and none was made; every reference inside the delivered unit resolves inside it, so no unresolved reference or citation is produced. Printed slips kept literal: no printed word, symbol, hypothesis or number of the counted statement or of its proof was altered. Layout and class: the article's own class is \texttt{amsart} with packages including euscript, comment, xcolor, amsthm, amsmath, amsfonts; the wrapper is the lane's pinned \texttt{amsart} editorial wrapper at 10pt on A4, which restates the theorem-like environments the retained text needs and adds the article's own macro definitions that the retained text needs (\texttt{\textbackslash codim}), with extra wrapper packages (bm, bbm, dsfont, xspace, cancel, mathtools). The delivered numbering is the unit's own. Assets: no retained span contains a figure environment, a graphics-inclusion command, a PDF-page inclusion, a table or a TikZ picture; no image file is copied into this package, the package declares no asset dependency and no image file is redistributed with it. Licence: version arXiv:2508.04320v1 of this article is distributed under the Creative Commons Attribution 4.0 International licence, as recorded in the arXiv licence metadata for this exact version and shown on the abstract page https://arxiv.org/abs/2508.04320v1. A full rights notice is delivered at licenses/arxiv-2508.04320-CC-BY-4.0.txt. Characters: every retained excerpt is pure ASCII, so the delivered engine, XeLaTeX with its default Latin Modern text font, reports no missing character.
\begin{equation}\label{LDE}
x(n+1)=\begin{pmatrix}
A_{11}(n) & A_{12}(n)\\
0 & A_{22}(n)
\end{pmatrix} x(n), \quad n\in J'.
\end{equation}

\begin{equation}\label{ldedd}
    x(n+1)=\begin{pmatrix}
A_{11}(n) & 0\\
0 & A_{22}(n)
    \end{pmatrix} x(n), \quad n\in J'.
\end{equation}

\begin{equation}\label{iid}
x_i(n+1)=A_{ii}(n)x_i(n) \quad n\in J'
\end{equation}

\begin{equation}\label{Si}
\mathcal S_i=\left \{v_i\in X_i: \ \sup_{n\in J}\|\Phi_{ii}(n, 0)v_i\|<+\infty \right \}, \quad i=1, 2.
\end{equation}

\begin{theorem}\label{thm7}
Let $J=\mathbb Z^+$ and suppose that~\eqref{LDE} admits an exponential dichotomy with projections $P(n)$ such that $d:=\dim \mathcal N P(0)<+\infty$. Furthermore, let $\mathcal S_i$, $i=1, 2$ be given by~\eqref{Si}.  Then, the following holds:
\begin{enumerate}
    \item  $\mathcal S_2$ is complemented and closed,  $d_2:=\codim \mathcal S_2 \le d$ and~\eqref{iid} admits an exponential dichotomy for $i=2$;
    \item $\mathcal S_1$ given by~\eqref{Si} for $i=1$ is closed and complemented;
    \item \eqref{ldedd} admits an exponential dichotomy if and only if
    \begin{equation}\label{dddeq}d=d_1+d_2,\end{equation} where $d_1:=\codim \mathcal S_1$.
    \end{enumerate} 
\end{theorem}

\begin{lemma}\label{ddd}
We have $d\ge d_1+d_2$.
\end{lemma}

\begin{theorem}\label{515thm}
Let $J=\mathbb Z^+$ and assume that $A_{ii}(n)$ is an invertible operator for $i=1, 2$ and $n\in J$. Suppose that~\eqref{LDE} admits an exponential dichotomy with respect to projections $P(n)$, $n\in J$, and that $d:=\dim \mathcal NP(0)<+\infty$. Then \eqref{ldedd} admits an exponential dichotomy.

\end{theorem}

\begin{proof}
It follows from Theorem~\ref{thm7} that~\eqref{iid} admits an exponential dichotomy for $i=2$. Moreover, $\mathcal S_1$ (given by~\eqref{Si} for $i=1)$ is closed. 
Next, we have
\[
\mathcal R P(0)=\{(\xi+L\eta, \eta): \xi \in \mathcal S_1, \ \eta \in \mathcal S_2'\},
\]
where $\mathcal S_2'$ is as in the proof of Lemma~\ref{ddd} and $L\colon \mathcal S_2'\to \mathcal S_1^\perp$ is a linear operator given by $L\eta=x_1(0)$, where $(x_1(n))_{n\in J}\subset X_1$ is the unique bounded solution of the equation
\[
x_1(n+1)=A_{11}(n)x_1(n)+A_{12}(n)\Phi_{22}(n, 0)\eta \quad n\in J,
\]
with $x_1(0)\in \mathcal S_1^\perp$. Note that 
\[
X_1\times X_2=\mathcal RP(0)\oplus (\mathcal S_1^\perp \times W),
\]
where $W$ is any subspace of $X_2$ such that $X_2=\mathcal S_2'\oplus W$. Hence, $\dim \mathcal S_1^\perp<+\infty$.
\begin{lemma}\label{bounded}
$L$ is a bounded linear operator.
\end{lemma}

\begin{proof}[Proof of the lemma]
Since $\dim \mathcal S_1^\perp<+\infty$, it suffices to show that $\mathcal N L$ is closed. To this end, we take a sequence $(\eta_k)_k \in \mathcal N L$ that converges to $\eta$ in $X_2$. Then $(0, \eta_k)_k$ is a sequence in $\mathcal RP(0)$ that converges to $(0, \eta)$. Since $\mathcal RP(0)$ is closed, we have $(0, \eta)\in \mathcal R P(0)$. In particular, $\eta \in \mathcal S_2'$ and $L\eta=0$
Hence, $\eta \in \mathcal N L$.
\end{proof}
\begin{lemma}\label{closed}
$\mathcal S_2'$ is closed.
\end{lemma}

\begin{proof}[Proof of the lemma]
Take a sequence $(\eta_k)_k$ in $\mathcal S_2'$ that converges to $\eta\in X_2$. Since $\mathcal S_2$ is closed, we have $\eta \in \mathcal S_2$.

By $(x_1^k(n))_{n\in J}$ we will denote the bounded sequence in $X_1$ such that 
\begin{equation}\label{x11k}
    x_1^k(n+1)=A_{11}(n)x_1^k(n)+A_{12}(n)\Phi_{22}(n, 0)\eta_k \quad n\in J,
\end{equation}
and $x_1^k(0)=L\eta_k \in \mathcal S_1^\perp$. In particular, for each $k\in \mathbb N$,
\[
(x_1^k(n), \Phi_{22}(n, 0)\eta_k)=\Phi(n, 0)(L\eta_k, \eta_k), \quad n\in J.
\]
Since $(L\eta_k, \eta_k)\in \mathcal R P(0)$, we have
\[
\|x_1^k(n)\| \le \|\Phi(n, 0)(L\eta_k, \eta_k)\| \le Ce^{-\lambda n} \|(L\eta_k, \eta_k)\| \le C(1+\|L\|)\sup_{k\in \mathbb N}\|\eta_k\|,
\]
for $n\in J$ and $k\in \mathbb N$, where $C, \lambda>0$ are constants associated with the exponential dichotomy of~\eqref{LDE}. Consequently, 
\[
\|x_1^k(n)\| \le D \quad \text{for $n\in J$ and $k\in \mathbb N$,}
\]
where $D:=C(1+\|L\|)\sup_{k\in \mathbb N}\|\eta_k\|$.

Since $X_1$ is reflexive (as it is a Hilbert space), for each $n\in J$, the sequence $(x_1^k(n))_k$ has a weakly convergent subsequence whose limit belongs to the closed ball in $X_1$ with radius $D$ centered in $0$. Using the diagonal procedure, we can find a subsequence $(k_j)_j$ of $\mathbb N$ such that $(x_1^{k_j}(n))_j$ converges weakly to $x_1(n)\in X_1$ with $\|x_1(n)\| \le D$, for every $n\in J$.
Applying~\eqref{x11k} for $k=k_j$, passing to the limit when $j\to \infty$, and recalling that weak convergence is preserved under the action of a bounded linear operator, we get
\[
x_1(n+1)=A_{11}(n)x_1(n)+A_{12}(n)\Phi_{22}(n, 0)\eta \quad n\in J.
\]
Therefore, $\eta\in \mathcal S_2'$. 
\end{proof}
Let $P_{11} \colon X_1\to \mathcal S_1$ be the projection on $\mathcal S_1$ along $\mathcal S_1^\perp$, and $P_{22}\colon X_2 \to \mathcal S_2'$ be the projection on $\mathcal S_2'$ along $(\mathcal S_2')^\perp$. We define the projection $\tilde P$ on $X_1\times X_2$ by 
\begin{equation}\label{tildeP}
\tilde P=\begin{pmatrix}
P_{11} & LP_{22}\\
0 &P_{22}
\end{pmatrix}.
\end{equation}
From the previous two lemmas it follows that $\tilde P$ is bounded. Moreover, $\mathcal R\tilde P=\mathcal RP(0)$. Set
\[
\tilde P(n):= \Phi(n, 0)\tilde P\Phi(0, n), \quad n\in J.
\]
Note that 
\[
\tilde P(n)=\begin{pmatrix}
\Phi_{11}(n, 0)P_{11}\Phi_{11}(0, n) & \star\\
0 &\Phi_{22}(n, 0)P_{22}\Phi_{22}(0, n)
\end{pmatrix}, \quad n\in J.
\]
Then \eqref{LDE} admits an exponential dichotomy with respect to the sequence of projections $\tilde P(n)$, $n\in J$. This easily implies that~\eqref{iid} for $i=2$ admits an exponential dichotomy with respect to the sequence of projections $\tilde P_{22}(n)$, $n\in J$ given by 
\[
\tilde P_{22}(n)=\Phi_{22}(n, 0)P_{22}\Phi_{22}(0, n), \quad n\in J.\]
%Consequently, $\mathcal S_2'=\mathcal R P_{22}=\mathcal S_2$, which implies that~\eqref{iid} for $i=2$ admits an exponential dichotomy.
\end{proof}

\end{document}
